\section{Introdução}
\textbf{Objetivos:}
\begin{itemize}
    \item Implementar os métodos de Newton, Newton Modificado e Newton Discreto.
    \item Analisar as regiões de convergência em sistemas de equações não lineares.
\end{itemize}
\vspace{5pt}
Seja
$F:\mathbb{R}^2 \to \mathbb{R}^2$ e $F(\mathbf{x}) = \mathbf{0}$ 
o sistema não linear a ser analisado.
Para cada solução 
$(x_r^*, y_r^*)$, com $r = 1, 2, \ldots, L$,
desse sistema, associamos uma cor diferente 
$C_r$ 
e seja 
$C_0$ 
uma cor distinta de todas as outras.

Seja 
$k_{\max}$ 
o número máximo de iterações permitidas pelo método de Newton e
$\epsilon > 0$ 
uma tolerância.

Considere uma região do plano discreta,
$T = \{(x_i, y_j) \mid i, j = 0, 1, \ldots, N\}$.
Para cada 
$(x_i, y_j) \in T$ 
vamos aplicar o método de Newton, a partir dessa aproximação inicial 
$\mathbf{x}^{(0)} = (x_i, y_j)$. 
Vamos definir uma função 
$C(x_i, y_j) = C_r$ se $\left\| (x^{(k)}, y^{(k)}) - (x_r^*, y_r^*) \right\| < \epsilon$
, sendo $r = 1, ... L$ a quantidade de soluções do sistema, e $C(x_i, y_j) = C_0$  caso contrário. Como resultado teremos uma aproximação da região de convergência em $T$ .
Este primeiro mapa ilustra para onde a solução aproximada convergiu.
Podemos também construir um segundo mapa de convergência que vai ilustrar como a sequência convergiu.
Nesse caso, a cor associada a cada $x^{(0)} = (x_i, y_j)$ aproximação inicial, está relacionada ao número de iterações
realizadas. Assim, vamos definir a função $M(x_i, y_j) = C_k$, sendo $C_k = 1, ..., k_{max}$ o número de iterações, e $C(x_i, y_j) = C_0$, caso contrário.
